\chapter{LFM2.5-24B-A2B}\label{ch:lfm25}

A decode token of LiquidAI's LFM2.5-24B-A2B reads 4.91~GB of its BF16 weights and 1.31~GB of our 4.0-bit EXL3 pack (\arith, from the headers), because only 4 of its 64 experts run in each of 38 layers. On one DGX Spark the engine decodes it at 136.8~tok/s from the 4.0 pack and 44.7~tok/s in BF16 (\measured, 2026-10-14). It is the third model TandemLLM serves and the first that is not a Qwen. This chapter says what is different about it and what is measured; Chapter~\ref{ch:adding-model} describes the adapter path it came in through, and \code{docs/lfm2.5-24b-a2b.md} is the same material for the repository.

\section{The model}\index{LFM2.5}

\begin{table}[htbp]
\centering\small
\caption{LFM2.5-24B-A2B beside Flash-Next.}\label{tab:lfm-model}
\begin{tabularx}{\linewidth}{@{}L{.2\linewidth}YL{.3\linewidth}@{}}
\toprule
\hd{part} & \hd{LFM2.5-24B-A2B} & \hd{Flash-Next} \\
\midrule
layers & 40: 30 short convolutions, 10 attention & 48: 36 Gated DeltaNet, 12 attention \\
MLP & 64 experts a layer, 4 run a token, no shared expert; layers 0 and 1 dense & 512 experts, 10 routed and 1 shared \\
residual stream & one of 2,048 & 4 streams of 2,560 \\
attention & full, 32 query and 8 KV heads of 64 & sparse past 2,051 tokens \\
state of a sequence & 240~KiB of convolution inputs plus 20,480 bytes a token & 110.3~MiB plus 29,952 bytes a token \\
context & 131,072 tokens & 262,144 \\
draft layer & none & one MTP layer \\
weights served & BF16, or our EXL3 packs at 2.5 to 6.0 bits a weight & turboderp's EXL3 packs \\
\bottomrule
\end{tabularx}
\end{table}

A short-convolution layer projects the normed input to three vectors B, C and x of 2,048, takes B times x, runs a depthwise convolution of 3 taps over that product along the sequence, multiplies by C and projects back. Its state is the last two products, so a sequence keeps 30 layers times 2 times 2,048 values. The router takes the sigmoid of each expert's logit, adds a per-expert bias to choose the 4, and renormalises the 4 chosen sigmoids. \code{engine/models/lfm2_moe/config.py} reads the configuration and stops on a field it does not know.

The tokenizer is LiquidAI's own (128,000 ids). The lookup corpus and the suffix store of the Qwen models hold other token ids, so the server refuses them by fingerprint and the LFM2.5 profiles keep a store of their own.

\section{Profiles}

\begin{table}[htbp]
\centering\small
\caption[The LFM2.5 profiles]{The LFM2.5 profiles. Weights: the GiB on the device at load, \measured{} for BF16 and 4.0; the other packs are estimated from the parameter count until each is measured.}\label{tab:lfm-profiles}
\begin{tabularx}{\linewidth}{@{}L{.22\linewidth}L{.12\linewidth}L{.14\linewidth}Y@{}}
\toprule
\hd{profile} & \hd{disk} & \hd{weights} & \hd{note} \\
\midrule
\code{lfm25-bf16} & 47.9~GB & 44.95~GiB & the quality reference \\
\code{lfm25-exl3-4.0} & 12.64~GB & 12.05~GiB & the default, provisional \\
\code{lfm25-exl3-2.5} to \code{-6.0} & 8.1 to 18.3~GB & & not measured yet \\
\bottomrule
\end{tabularx}
\end{table}

Every profile fits the Spark preset at the model's whole context: its KV is 2.5~GiB at 131,072 tokens. The default, \code{lfm25-exl3-4.0}, was chosen before any pack was measured, as the smallest whose rates are all whole bits, and the measured decode rate with a chat-quality floor decides it. On 8 chat prompts of 128 tokens, the 4.0 pack's top token agrees with transformers' BF16 model at 95.6~\% of the 1,024 positions, with a KL divergence of 0.027 (\measured, the engine teacher-forced on the BF16 run's tokens).

Below 4 bits a weight, exllamav3 1.6.0 gives some tensors half-integer rates (1.5, 2.5 or 3.5 bits; a tile is $16K + 8$ int16 wide). The engine's EXL3 kernels decode whole-bit rates, so the pack reader stops on such a tensor and names it.

\section{One row path and exact rollback}

The decode step and every row of a verify block run the same kernels a row at a time (\code{engine/models/lfm2_moe/kernels.py}): the norms, the short convolution over the row's own window, q and k with their rotary, attention over the row's visible keys in fixed chunks of 512 keys, and the router. The EXL3 projections run on \code{tools/fn_kernels.dense} and the experts on its MoE bank kernels, which Flash-Next uses too; the BF16 projections run on a 16-row tile GEMM. So a row of a verify block has the bits of the decode step at that position. The load-time self-check proves it on the device: a recorded decode of 24 tokens verified as one block, every row bit for bit, before speculation turns on.

A verify block writes nothing into the convolution state. The commit takes the new state from the last two inputs along the accepted path, which are copies of the block's own rows, and moves the path's KV rows down. A rollback is exact by construction.

\begin{table}[htbp]
\centering\small
\caption[The LFM2.5 correctness ladder]{The correctness ladder. \measured, 2026-10-14.}\label{tab:lfm-ladder}
\begin{tabularx}{\linewidth}{@{}L{.36\linewidth}Y@{}}
\toprule
\hd{check} & \hd{result} \\
\midrule
reference forward against transformers, a tiny random model in fp32 & within 2.1e-5 over a prefill and 3 decode steps \\
BF16 engine against transformers in BF16, 8 chat prompts, 1,024 positions & top token 98.05~\%; 19 of the 20 differences near ties; KL 0.0064 \\
4.0 engine against exllamav3's own forward, the same tokens & top token 97.85~\%; all 22 differences near ties \\
trellis decode against exllamav3's \code{LinearEXL3} & bit for bit \\
\bottomrule
\end{tabularx}
\end{table}

\section{Speed}

One request, greedy, the decode step from its CUDA graph: 136.8~tok/s on 4.0 (7.31~ms a token, 179~GB/s of the bytes a token reads) and 44.7~tok/s in BF16 (22.4~ms, 220~GB/s). Launching the step's kernels one by one costs 5~\% (129.8 and 42.7~tok/s). An 8,192-token prompt prefills at 7,360~tok/s on 4.0 (\measured).

A verify block costs about 3~ms a row more, because each row routes to its own experts and a block reads the union of them (Table~\ref{tab:lfm-stair}). A block pays only when it is accepted for about 1.4 tokens a row it adds.

\begin{table}[htbp]
\centering\small
\caption[The verify staircase on 4.0]{A chain block of $n$ rows on 4.0 from its graph, the median of 20 replays at 1k tokens of context, and the experts its rows read over the 38 MoE layers (of 2,432). \measured, 2026-10-14.}\label{tab:lfm-stair}
\begin{tabular}{@{}lrrrrrrrr@{}}
\toprule
\hd{rows} & 1 & 2 & 3 & 4 & 6 & 8 & 12 & 16 \\
\midrule
ms & 7.42 & 10.62 & 13.79 & 16.86 & 21.40 & 27.47 & 36.17 & 45.88 \\
experts & 152 & 259 & 333 & 391 & 491 & 563 & 664 & 735 \\
\bottomrule
\end{tabular}
\end{table}

The model has no draft layer, and Liquid will publish a drafter for it. Until then the server drafts with the lookup alone, through the shared tree drafter with no head (\code{NoHead}) and cut by StairCut with the staircase as its prior. Through the server on 4.0 the lookup gave the same answer as plain decoding for all 20 requests of \code{tools/fn_server_check.py} (greedy, thinking, streaming, every \code{tool_choice}, JSON, penalties with \code{logit_bias}, logprobs, seeds, a stop string), and the largest logprob difference was 0.0. Those prompts took 1.0 to 1.2 tokens a block, and the speed did not change.

\section{Serving}

The chat template is the checkpoint's own and opens a reasoning block at the end of every prompt; a client that turns thinking off gets the block closed right after the generation prompt. Tool calls are a Python list between \code{<|tool_call_start|>} and \code{<|tool_call_end|>}; \code{server/toolcall.py} rewrites each block into the parser's JSON form as the text arrives, with values read as Python literals and never run, so a client gets OpenAI \code{tool_calls}. Liquid's recommended sampling is the default: temperature 0.1, \code{top_k} 50, repetition penalty 1.05.

The resident prefix and the state store keep the state through the adapter. A warm server's answers, with the prefix resumed at 4,096 and 12,288 tokens, equalled a cold server's with equal logprobs (\code{tools/fn_cache_check.py}, \measured). Side slots run over the one copy of the weights, and a slot's plain step joins the other slots' round as one row. With 2 and 4 requests at once every round was a joint pass and every answer equalled its request alone; that pass ran eagerly and the total stayed at 126 to 130~tok/s.

\section{What is not measured}

\begin{itemize}
\item The packs other than 4.0, and which profile becomes the default.
\item The joint decode step from its graph with several slots.
\item The disk tier and the session spill on this model; the profiles turn them off.
\item Half-integer EXL3 rates and 7-bit experts, which the kernels do not decode yet.
\item Any task benchmark, and the benchmark row of Chapter~\ref{ch:measurement}.
\end{itemize}
